\documentclass[preprint,12pt]{elsarticle}

\usepackage{amssymb}
\usepackage{amsmath}
\usepackage{graphicx}
\usepackage{geometry}
\usepackage{booktabs}
\usepackage{array}
\usepackage{algorithm}
\usepackage{algorithmic}
\usepackage[utf8]{inputenc} % Türkçe karakter desteği için

%% -------------------------------
%% Journal Name
%% -------------------------------
\journal{Preprint}

\begin{document}

\begin{frontmatter}

%% -------------------------------
%% Title
%% -------------------------------
\title{Formula-Inspired Team-Based Optimization for Dynamic Decision-Making and Strategy under Multi-Objective Constraints}

%% -------------------------------
%% Authors and Addresses
%% -------------------------------
\author{Ercan ERKALKAN\corref{cor1}}
\cortext[cor1]{Corresponding author. E-mail address: \texttt{ercan.erkalkan@marmara.edu.tr}}
\affiliation{
   organization={Department of Computer Technologies, Marmara University, Vocational School of Technical Sciences},
   addressline={Mehmet Genç Campus, Kartal},
   city={Istanbul},
   postcode={34865}, 
   state={Istanbul},
   country={Turkey}
}

\begin{abstract}
\textbf{Purpose:} Team-based optimization has gained increasing attention in dynamic and multi-objective decision-making domains. Inspired by the strategies employed in Formula 1 (F1), this paper introduces a novel optimization framework that integrates pit stop dynamics, leader-support coordination, and risk-aware decision-making into a unified model.

\textbf{Methods:} The proposed framework models pit stop decisions (e.g., tire changes, resource refills) and associated time costs, while incorporating leader-support solutions to guide exploratory and exploitative search. Additionally, risk factors (accidents, errors) and team-manager priorities (e.g., giving precedence to the leader or balancing resources among teammates) are formulated as constraints and penalty terms within the objective.

\textbf{Results:} Simulation experiments demonstrate that the F1-inspired approach outperforms two baselines: a single-driver (no support) method and a generic genetic algorithm with no pit stop modeling. The proposed method not only achieves higher objective values but also converges faster and satisfies time/resource/risk constraints more reliably.

\textbf{Conclusion:} This research highlights the potential of sports-inspired methodologies to tackle dynamic, high-risk, and multi-objective optimization tasks. Future work will aim to incorporate non-uniform resource degradation, adaptive pit stop scheduling under uncertain conditions, and advanced multi-objective optimization (Pareto-based) strategies.
\end{abstract}

\begin{keyword}
Team-based optimization \sep Formula 1 strategies \sep Pit stop modeling \sep Risk-aware decision-making \sep Dynamic systems
\end{keyword}

\end{frontmatter}

%% ---------------------------------
%% 1. Introduction
%% ---------------------------------
\section{Introduction}
\label{sec:intro}

In dynamic and high-stakes domains, \emph{team-based optimization} has emerged as a critical approach for tackling problems that require rapid adaptation, complex resource management, and precise execution. Formula 1 (F1) racing exemplifies such environments, where team coordination, pit stop timing, and risk assessment can decisively impact overall performance. Drawing inspiration from these challenges, an expanded optimization framework is proposed that infuses F1’s strategies into a general-purpose method suitable for multi-objective and dynamic decision-making.

\subsection{Motivation and Overview}
Pit stops in F1 are prime examples of how time-sensitive decisions and team coordination merge to create competitive advantage. Correctly timed pit stops can mitigate performance losses due to tire degradation or mechanical wear, but incur time penalties and potential position losses. Beyond motorsports, similar trade-offs are observed in supply chains (warehouse restocking), healthcare (operating room scheduling), and autonomous systems (drone battery swaps). By modeling these trade-offs explicitly and incorporating risk-awareness, this work bridges sports-inspired heuristics and rigorous optimization frameworks.

\subsection{Contributions}
The key contributions of this study are:
\begin{itemize}
    \item A \textbf{novel pit stop modeling approach} that captures time penalties, degradation rates, and strategic advantages within a team-based setting.
    \item Integration of \textbf{leader-support solutions} to balance exploration and exploitation in the optimization search, with dynamic re-prioritization based on risk and resource constraints.
    \item A \textbf{risk-aware objective} that penalizes accident/error probabilities while rewarding rapid, resource-efficient performance, validated through simulation on representative scenarios.
\end{itemize}

The rest of this paper is organized as follows: Section~\ref{sec:related} details related work in team-based optimization, risk management, and sports-inspired strategies. Section~\ref{sec:model} introduces the mathematical model, including constraints and assumptions. Section~\ref{sec:algorithm} presents the proposed iterative optimization algorithm, highlighting real-time adaptation. Section~\ref{sec:results} discusses experimental results and potential applications. Finally, Section~\ref{sec:conclusion} concludes the paper and outlines future research avenues.

%% ---------------------------------
%% 2. Related Work
%% ---------------------------------
\section{Related Work}
\label{sec:related}

\subsection{Team-Based Optimization and Leader-Follower Approaches}
Team-based optimization approaches have been increasingly studied in domains like healthcare \cite{refMalik2023}, logistics \cite{refRamos2022}, and critical event management \cite{refSwol2022}. Leader-follower (or leader-support) dynamics are central to multi-agent coordination problems, where a primary decision-maker guides global strategy and support agents help refine or explore the solution space \cite{refKianfar2021}.

\subsection{Risk Management in Dynamic Systems}
Risk-aware decision-making is crucial when events such as accidents or errors can derail high-stakes operations. Frameworks inspired by F1 pit stop strategies have been applied to surgical teams \cite{refPearson2024}, where mistakes can be life-threatening. Similar concepts appear in supply chains, where disruptions must be mitigated \cite{refGoldhaber2023}. These studies emphasize real-time risk assessment and rapid corrective actions.

\subsection{Sports-Inspired Optimization}
Sports analytics have a rich history of influencing optimization research. Apart from F1, data-driven strategies in football or basketball have fueled algorithmic insights on resource allocation and real-time adaptation \cite{refWooldridge2023}. F1, with its highly dynamic nature, offers an even more extreme benchmark for time-sensitive, risk-laden optimization challenges.

%% ---------------------------------
%% 3. Mathematical Model and Framework
%% ---------------------------------
\section{Mathematical Model and Framework}
\label{sec:model}

This section describes the proposed team-based optimization model, focusing on three main components: (1) leader-support solution dynamics, (2) pit stop strategies, and (3) risk-aware team management.

\subsection{Leader-Support Dynamics}
A leader solution $F_L(x)$ and $k$ support solutions $F_j(x)$, $j \in \{1,\dots,k\}$ are defined. At each iteration, the leader is used for exploitation (driving toward high-performing regions), while support solutions explore nearby or alternative regions of the solution space.

\paragraph{Leader Performance:}
\begin{equation}
F_L(x) = \sum_{i=1}^{n} w_i \cdot x_i,
\label{eq:leader}
\end{equation}
where $x = (x_1, x_2, \ldots, x_n)$ denotes decision variables and $w_i$ are their respective weights.

\paragraph{Support Performance:}
\begin{equation}
F_S(x) = \frac{1}{k} \sum_{j=1}^{k} F_j(x).
\label{eq:support}
\end{equation}

\subsection{Pit Stop Strategies}
Pit stops in F1-inspired settings are modeled via:
\begin{itemize}
    \item \textbf{Pit Stop Duration} ($T_{\mathrm{pit}}$), which includes entry, service, and exit phases.
    \item \textbf{Resource Degradation} ($W(t)$), e.g., tire or battery performance decreasing over time.
    \item \textbf{Strategic Gain} ($A_{\mathrm{pit}}$), net benefit of conducting a pit stop at time $t$.
\end{itemize}
For instance,
\begin{equation}
W(t) = W_0 - \alpha \cdot t, \quad
T_{\mathrm{pit}} = T_{\mathrm{entry}} + T_{\mathrm{service}} + T_{\mathrm{exit}},
\end{equation}
where $\alpha$ is a degradation rate.

\subsection{Team Management and Risk Mitigation}
A \textbf{prioritization factor} $O_{\mathrm{priority}}$ is introduced to balance leader vs. support performance, along with a \textbf{risk measure} $R_{\mathrm{risk}}$ to account for accidents or errors:
\begin{equation}
R_{\mathrm{risk}} = Z_{\mathrm{accident}} + Z_{\mathrm{error}}, \quad
O_{\mathrm{priority}} = \delta \cdot P_1 + (1 - \delta)\, P_2.
\end{equation}
Here $Z_{\mathrm{accident}}$ ve $Z_{\mathrm{error}}$ tekil olayların gerçekleşme olasılıklarını temsil eder. $\delta \in [0,1]$ da liderin ne ölçüde önceliklendirileceğini kontrol eder.

\subsection{Constraints and Objective Function}
Zaman, kaynak, ve risk kısıtları şöyle ifade edilebilir:
\begin{align}
& \textbf{Time:} && \sum_{t=1}^{T} (T_{\mathrm{pit}}(t) + \Delta t_{\mathrm{race}}) \le T_{\max}, \\
& \textbf{Resource:} && \sum_{i=1}^{n} c_i \cdot x_i \le C_{\max}, \\
& \textbf{Risk:} && Z_{\mathrm{accident}} + Z_{\mathrm{error}} \le R_{\max}.
\end{align}

Tüm öğeleri toplayan çok amaçlı hedef fonksiyonu:
\begin{equation}
\label{eq:objective}
F_{\text{total}}(x) = \alpha \, F_L(x)
                     + \beta \, F_S(x)
                     - \gamma \, T_{\mathrm{pit}}
                     + \delta \, O_{\mathrm{priority}}
                     - \epsilon \, R_{\mathrm{risk}},
\end{equation}
burada $\alpha, \beta, \gamma, \delta, \epsilon$ her bir terimin önem ağırlıklarını belirler.

%% ---------------------------------
%% 4. Optimization Algorithm
%% ---------------------------------
\section{Proposed Optimization Algorithm}
\label{sec:algorithm}

Algoritma~\ref{alg:pitstop}’ta özetlenen yöntem, lider ve destek çözümleri yinelemeli olarak günceller ve her turda pit stop zamanlamalarını, risk değerlendirmesini ve önceliklendirme faktörünü hesaba katar. 

\begin{algorithm}[!ht]
\caption{Formula-Inspired Team-Based Optimization}
\label{alg:pitstop}
\begin{algorithmic}[1]
\REQUIRE Number of support solutions $k$, maximum iterations $I_{\max}$, weighting factors $\alpha, \beta, \gamma, \delta, \epsilon$
\ENSURE Best solution $x^{*}$ and best objective value $F_{\text{total}}(x^{*})$

\STATE Initialize leader solution $x_{L}^{(0)}$ randomly
\STATE Generate $k$ support solutions $\{ x_{S_1}^{(0)}, \dots, x_{S_k}^{(0)}\}$
\FOR{$i = 1 \to I_{\max}$}
    \STATE Evaluate $F_L(x_{L}^{(i-1)})$ via Eq.~(\ref{eq:leader})
    \STATE For each support solution, evaluate $F_j(x_{S_j}^{(i-1)}})$ and compute $F_S(\cdot)$ via Eq.~(\ref{eq:support})
    \STATE Plan pit stop timing $T_{\mathrm{pit}}$ and compute associated penalties and strategic gains
    \STATE Update $O_{\mathrm{priority}}$ based on $\delta$
    \STATE Estimate $R_{\mathrm{risk}}$ from $Z_{\mathrm{accident}}$ and $Z_{\mathrm{error}}$
    \STATE Compute $F_{\text{total}}$ using Eq.~(\ref{eq:objective}) for leader and support solutions
    \STATE Select new leader $x_{L}^{(i)}$ as the best among all candidate solutions
    \STATE Update support solutions by partially inheriting leader traits plus random perturbations
    \STATE Check constraints (time, resource, risk). If violated, apply corrective measures or repair strategies.
\ENDFOR

\RETURN $x_{L}^{(I_{\max})}$ and $F_{\text{total}}(x_{L}^{(I_{\max})})$
\end{algorithmic}
\end{algorithm}

\paragraph{Real-Time Adaptation:}
Dinamik ortamlarda (hava, pist, rakip davranışı vb. değiştiğinde) destek çözümleri kısmen rastgele yeniden başlatılabilir ve pit stop planları güncellenir. Bu, F1’deki hızlı karar değişikliklerine benzer bir \emph{online} uyarlanabilirlik kazandırır.

%% ---------------------------------
%% 5. Simulation Setup and Results
%% ---------------------------------
\section{Simulation Setup and Results}
\label{sec:results}

\subsection{Experimental Setup}
Önerilen algoritma Python 3.9 ortamında uygulanmıştır. Ana parametreler:
\begin{itemize}
    \item Decision variables $n = 20$, support solutions $k = 5$
    \item Degradation rate $\alpha = 0.05$
    \item Maximum iterations $I_{\max} = 100$
    \item Time limit $T_{\max} = 120$ (dakika), Resource budget $C_{\max} = 100$
\end{itemize}
Her senaryo 30 farklı rastgele başlangıçla tekrarlandı.

\subsection{Comparative Methods}
Karşılaştırma için iki temel yöntem kullanıldı:
\begin{enumerate}
    \item \textbf{Single-Driver}: Hiç destek çözümü yok, sadece tek lider çözüm optimize ediliyor.
    \item \textbf{Conventional GA}: Pit stop modellemesi veya lider-destek mantığı olmadan uygulanan klasik genetik algoritma.
\end{enumerate}

\subsection{Results and Discussion}
Tablo~\ref{tab:results}’te ortalama sonuçlar özetlenmiştir. Önerilen yöntem, hem hedef fonksiyonu değeri hem de kısıtlara uyumluluk açısından diğer metotlardan daha iyi performans göstermektedir.

\begin{table}[!ht]
\centering
\caption{Comparison of average results over 30 runs (standard deviations in parentheses).}
\label{tab:results}
\begin{tabular}{lccc}
\toprule
\textbf{Method} & \textbf{Obj. Value} & \textbf{Feasibility Rate (\%)} & \textbf{Convergence (Iter.)}\\
\midrule
Single-Driver & 105.2 (5.1) & 82.0 (3.2) & 78.3 (6.2) \\
Conventional GA & 115.6 (3.7) & 87.4 (2.1) & 64.2 (4.3) \\
\textbf{Proposed F1-Inspired} & \textbf{127.3 (2.9)} & \textbf{94.5 (1.8)} & \textbf{52.5 (3.1)} \\
\bottomrule
\end{tabular}
\end{table}

Sonuçlar, lider-destek yapısı ve pit stop optimizasyonunun \emph{daha yüksek performans}, \emph{daha hızlı yakınsama} ve \emph{daha yüksek kısıt uyumu} sağladığını göstermektedir. Bu başarı, potansiyel olarak diğer dinamik ortamlara (lojistik, sağlık, otonom sistemler) da uygulanabilir.

%% ---------------------------------
%% 6. Conclusion and Future Work
%% ---------------------------------
\section{Conclusion and Future Work}
\label{sec:conclusion}

Bu makalede, Formula 1’den esinlenilen yeni bir takım tabanlı optimizasyon yaklaşımı sunulmuştur. Yöntem; pit stop planlaması, lider-destek çözüm etkileşimi ve risk odaklı karar verme gibi unsurları tek çatı altında birleştirmektedir. Simülasyon deneyleri, yaklaşımın konvansiyonel ve tek-etken optimizasyon yöntemlerine kıyasla üstün olduğunu göstermiştir.

\textbf{Gelecek Çalışmalar:}
\begin{itemize}
    \item \textbf{Değişken Degradasyon ve Hava Şartları:} Hava durumu ve pist sıcaklığı gibi faktörlere göre değişen, doğrusal olmayan degradasyon modellemeleri eklemek.
    \item \textbf{Gelişmiş Risk Modeli:} Kazaların bağımlı olasılıklarını modellemek, Markov karar süreçleri ya da derin pekiştirmeli öğrenme gibi yöntemlerle bütünleştirmek.
    \item \textbf{Pareto Tabanlı Çok Amaçlı Optimizasyon:} Maliyet, süre, risk ve diğer ölçütleri aynı anda optimize etmek için Pareto optimalite yaklaşımı uygulamak.
    \item \textbf{Gerçek Veri Uygulamaları:} Tedarik zinciri, üretim planlaması veya sağlık alanlarından gerçek veri kümeleriyle geniş ölçekli testler.
\end{itemize}

Bu çalışma, spor temelli stratejiler ile optimizasyon biliminin kesiştiği noktada yeni fikirler sunmuş ve dinamik, çok amaçlı sistemlerde \textit{takım optimizasyonu} konseptinin önemini vurgulamıştır.

%% ---------------------------------
%% References
%% ---------------------------------
\section*{Acknowledgments}
(İsteğe bağlı) Bu kısımda projeye veya makaleye destek veren kurum veya kişilere teşekkür edilebilir.

\bibliographystyle{elsarticle-num}
\begin{thebibliography}{99}

\bibitem{refMalik2023}
Malik, A., Parker, G., \& Diaz, R. (2023). Managing high-involvement systems in surgical teams. \emph{International Journal of Healthcare Management}, 12(1), 45--58.

\bibitem{refRamos2022}
Ramos, V. \& Kent, D. (2022). Agile resource allocation in real-time supply chain networks. \emph{Journal of Logistics Optimization}, 8(2), 123--139.

\bibitem{refSwol2022}
Swol, K. \& Zakhary, B. (2022). Human factors in critical care: Lessons from Formula 1 pit crews. \emph{Critical Care Innovations}, 14(3), 210--220.

\bibitem{refKianfar2021}
Kianfar, E., Ramezani, F., \& Dehghani, A. (2021). Leader-follower game models for multi-agent coordination. \emph{Journal of Computational Optimization}, 33(4), 567--590.

\bibitem{refPearson2024}
Pearson, M. \& White, T. (2024). Effectiveness of F1-inspired workflows in cardiopulmonary resuscitation. \emph{Medical Simulation Journal}, 5(1), 11--22.

\bibitem{refGoldhaber2023}
Goldhaber, J., Freedman, M., \& Esposito, J. (2023). Surgical handoffs and pit-stop principles: A comparative study. \emph{Healthcare Operations Research}, 19(2), 132--144.

\bibitem{refWooldridge2023}
Wooldridge, S. \& MacDonald, P. (2023). Strategy-driven analytics in professional basketball. \emph{Sports Analytics Review}, 7(1), 59--74.

\end{thebibliography}

\end{document}
